\scanpage{130}
\mspage{0-ii}
$\Gamma_U(f_\alpha):\Gamma(U,\underline F_\alpha)\longrightarrow\Gamma(U,\underline G_\alpha)$ forment un système projectif d'homomorphismes, les diagrammes
\diagram{\begin{tikzcd}
\Gamma(U,\underline F_\alpha)\arrow[r]\arrow[d]&\Gamma(U,\underline G_\alpha)\arrow[d]\\
\Gamma(V,\underline F_\alpha)\arrow[r]&\Gamma(V,\underline G_\alpha)
\end{tikzcd}}
\figure{\begin{tikzpicture}[x=0.3cm,y=0.3cm]\draw (0,0)--(0.02,0.5)--(-0.12,1.2);\draw (2.2,0)--(2.2,0.7)--(1.7,1.1);\end{tikzpicture}}
étant en outre commutatifs pour $U\supset V$; on en conclut que si $f_U=\varprojlim\Gamma_U(f_\alpha)$, les homomorphismes $f_U$ définissent un homomorphisme de faisceaux $f:\varprojlim\underline F_\alpha\longrightarrow\varprojlim\underline G_\alpha$, que l'on désigne par $\varprojlim f_\alpha$.

On peut donc dire que $\varprojlim$ est un \uline{foncteur} de la catégorie des systèmes projectifs de faisceaux de groupes abéliens sur $X$ dans la catégorie des faisceaux de groupes abéliens sur $X$. Ce foncteur est \uline{exact à gauche}: autrement dit, si \deleted{\unsure{pour chaque} \illegible{3}} $(\underline F'_\alpha)$, $(\underline F_\alpha)$, $(\underline F''_\alpha)$ sont trois systèmes projectifs de faisceaux de groupes abéliens sur $X$, et si pour chaque $\alpha$, on a une suite exacte avec commutativité des diagrammes
\diagram{\begin{tikzcd}
0\arrow[r]&\underline F'_\beta\arrow[r,"f_\beta"]\arrow[d]&\underline F_\beta\arrow[r,"g_\beta"]\arrow[d]&\underline F''_\beta\arrow[d]\\
0\arrow[r]&\underline F'_\alpha\arrow[r,"f_\alpha"]&\underline F_\alpha\arrow[r,"g_\alpha"]&\underline F''_\alpha
\end{tikzcd}}
pour $\alpha\leq\beta$, alors en posant $f=\varprojlim f_\alpha$, $g=\varprojlim g_\alpha$, la suite
\diagram{\begin{tikzcd}
0\arrow[r]&\varprojlim\underline F'_\alpha\arrow[r]&\varprojlim\underline F_\alpha\arrow[r]&\varprojlim\underline F''_\alpha
\end{tikzcd}}
est exacte. Cela résulte aussitôt de ce que le foncteur $\Gamma$ est exact à gauche, et de la propriété correspondante pour le foncteur $\varprojlim$ dans la catégorie des systèmes projectifs de groupes abéliens.

Soit $\varphi$ une application continue de $X$ dans un espace topologique $Y$. Si $(\underline F_\alpha)$ est un système projectif de faisceaux de groupes abéliens sur $X$, on a
\[
\varphi_*(\varprojlim\underline F_\alpha)=\varprojlim\varphi_*(\underline F_\alpha).
\]
En effet, pour tout ouvert $V\subset Y$; on a par définition
\[
\begin{aligned}
\Gamma(V,\varprojlim\varphi_*(\underline F_\alpha))&=\varprojlim\Gamma(V,\varphi_*(\underline F_\alpha))=\varprojlim\Gamma(\varphi^{-1}(V),\underline F_\alpha)\\
&=\Gamma(\varphi^{-1}(V),\varprojlim\underline F_\alpha)=\Gamma(V,\varphi_*(\varprojlim\underline F_\alpha))
\end{aligned}
\]
\margin{129}
